\subsection*{Pertanyaan 1: Buatlah fungsi QuickSort!}

\begin{lstlisting}[basicstyle=\ttfamily\scriptsize, frame=single, backgroundcolor=\color{white}, numbers=none, keywordstyle=\color{black}, commentstyle=\color{black}, stringstyle=\color{black}, breaklines=true]
def quick_sort(arr, i, j):
    # basis: array 1 elemen atau kosong
    if i >= j:
        return
    # partisi dan dapatkan indeks pembagi
    k = partisi(arr, i, j)
    # rekursif urutkan bagian kiri dan kanan
    quick_sort(arr, i, k)
    quick_sort(arr, k + 1, j)

if __name__ == "__main__":
    data = [4, 12, 23, 9, 21, 1, 35, 2, 24]
    print("Data awal:", data)
    
    # MinMaks D&C
    mn, mx = minmaks(data, 0, len(data) - 1)
    print(f"\nMinMaks D&C -> Min: {mn}, Maks: {mx}")
    
    # Merge Sort
    arr_ms = data.copy()
    merge_sort(arr_ms, 0, len(arr_ms) - 1)
    print(f"Merge Sort -> {arr_ms}")
    
    # Quick Sort
    arr_qs = data.copy()
    quick_sort(arr_qs, 0, len(arr_qs) - 1)
    print(f"Quick Sort -> {arr_qs}")
\end{lstlisting}

\subsection*{Pertanyaan 2: Buatlah fungsi Partition!}

\begin{lstlisting}[basicstyle=\ttfamily\scriptsize, frame=single, backgroundcolor=\color{white}, numbers=none, keywordstyle=\color{black}, commentstyle=\color{black}, stringstyle=\color{black}, breaklines=true]
def partisi(arr, i, j):
    # pilih elemen tengah sebagai pivot
    pivot = arr[(i + j) // 2]
    p, q = i, j
    while True:
        # scan dari kiri sampai menemukan elemen >= pivot
        while arr[p] < pivot:
            p += 1
        # scan dari kanan sampai menemukan elemen <= pivot
        while arr[q] > pivot:
            q -= 1
        # jika belum bertemu, tukar dan geser indeks
        if p < q:
            arr[p], arr[q] = arr[q], arr[p]
            p += 1
            q -= 1
        else:
            break
    return q
\end{lstlisting}